
\chapter{ПОД/ФТ}\label{ch:aml}
ПОД/ФТ~--- противодействие легализации (отмыванию) доходов и~финансированию
терроризма; в~международной практике для близкого набора задач используется
аббревиатура AML (\textit{anti-money laundering}). Платёжный комплаенс по~ПОД/ФТ
опирается на~115-ФЗ и~подзаконные акты Банка России и~Росфинмониторинга.
Международные стандарты заданы FATF (\textit{Financial Action Task Force},
Группа разработки финансовых мер борьбы с~отмыванием денег)~\cite{fatf-recs},
а~для СНГ (Содружество Независимых Государств) ту же методологическую роль
играет ЕАГ~--- Евразийская группа по~противодействию легализации преступных
доходов и~финансированию терроризма~\cite{eag-overview}; подробнее
см.~раздел~\ref{sec:aml-fatf-eag}.

\section{115-ФЗ и~платёжный бизнес}\label{sec:aml-115fz}

Федеральный закон от~07.08.2001~№~115-ФЗ \enquote{О~противодействии легализации
  (отмыванию) доходов, полученных преступным путём, и~финансированию
терроризма}~--- основа контроля ПОД/ФТ в~России~\cite{fz-115}. \highlight{Закон
  требует от~оператора знать клиента, наблюдать операции и~при формальных
основаниях блокировать средства}.

Первая обязанность~--- идентификация. Закон различает полную и~упрощённую
процедуру, а~подзаконные акты Банка России задают состав сведений и~документов
по~клиенту, представителю, выгодоприобретателю и~бенефициарному
владельцу~\cite{cbr-499p}. Регистрационных данных не~всегда достаточно: где-то
хватает минимального набора, а~где-то нужны документы, внешние реестры
и~дополнительные вопросы. Объём проверки зависит от~типа клиента, операции
и~уровня риска; пороги и~частные исключения регулятор пересматривает быстрее,
чем выходит редакция книги, поэтому фиксировать их здесь смысла нет. Подробно
идентификация физлиц и~ТСП разобрана в~главе~\ref{ch:kyc-kyb}.

Вторая обязанность~--- мониторинг: оператор замечает операции, подпадающие
под~критерии обязательного контроля, и~отправляет сведения в~Росфинмониторинг.
Этим задача не~ограничивается: надо ловить и~те операции, которые формально ещё
не~описаны в~законе, но~выбиваются из~нормального профиля клиента по~поведению
или по~контексту. Для этого нужны регламент, поток данных, маршрут проверки
и~понятное решение по~результату; контроль по~одним лишь формальным критериям
либо ничего не~увидит, либо утонет в~ложных срабатываниях.

Обязательный контроль по~статье~6 115-ФЗ применяется к~конечному перечню
операций~\cite{fz-115}. По~сумме базовый порог 1\,000\,000~руб.\ работает для
операций с~наличными у~юрлиц (внесение и~снятие со~счёта), для покупки наличной
валюты физлицом и~для ряда смежных категорий; для сделок с~недвижимостью порог
поднят до~5\,000\,000~руб., для зачислений на~отдельные счета
по~гособоронзаказу~--- 600\,000~руб., для почтовых переводов~--- 100\,000~руб.
По~виду операция уходит в~обязательный контроль независимо от~суммы, когда
сторона находится в~перечне лиц, причастных к~экстремизму или терроризму, либо
в~смежных перечнях Росфинмониторинга. Факультативный контроль по~статье~7
обязывает оператора выявлять необычные и~подозрительные операции по~собственным
правилам внутреннего контроля; именно эту работу закрывает программа выявления
операций в~ПВК.

Сведения направляются в~Росфинмониторинг через личный кабинет~---
формализованным электронным сообщением (ФЭС). Сроки задают 115-ФЗ и~нормативные
акты Банка России
и~Росфинмониторинга~\cite{fz-115,cbr-860p,rosfinmonitoring-oes}.

ФЭС собирается из~платёжного журнала, не~через отдельный процесс. Состав
показателей и~форматы регулируются Приказом Росфинмониторинга от~08.02.2022~№~18
\enquote{О~представлении в~Росфинмониторинг информации по~115-ФЗ} в~действующей
редакции~\cite{rosfinmonitoring-oes}, а~требования к~ПВК, на~которые опирается
логика выявления, устанавливает 860-П~\cite{cbr-860p}. В~инфраструктуре сборка
ФЭС идёт по~цепочке: программа выявления формирует запись с~кодом сработавшего
критерия; программа документального фиксирования дополняет её реквизитами
клиента и~контрагента из~программы идентификации; сборщик ФЭС упаковывает
результат в~формализованное сообщение к~моменту отправки. Если поле клиента
в~платёжном журнале не~заполнено или не~нормализовано, ФЭС не~соберётся.
Качество этих полей проверяется в~платёжной системе заранее, не~в~момент
срабатывания тревоги.

В~международной практике ту же роль играют сообщения SAR (\textit{Suspicious
Activity Report}) и~STR (\textit{Suspicious Transaction Report}), направляемые
в~национальный финансово-разведывательный орган~--- например, FinCEN в~США.
По~функции они эквивалентны российскому ФЭС; по~составу полей и~регламенту
подачи прямого соответствия форм нет.

Третья обязанность~--- замораживание денежных средств и~иного имущества. Если
клиент или бенефициар попадает в~перечни лиц, связанных с~терроризмом или
экстремизмом, организация обязана применить меры по~замораживанию и~исполнить
предусмотренный законом порядок информирования~\cite{fz-115}. Технология
сопоставления имён и~эскалации разобрана в~главе~\ref{ch:sanctions}.

\subsection*{Правила внутреннего контроля по~860-П}

Идентификацию, мониторинг и~замораживание банк собирает в~единый рабочий
регламент~--- правила внутреннего контроля (ПВК). Положение Банка России
от~18.06.2025~№~860-П \enquote{О~требованиях к~ПВК по~ПОД/ФТ}\footnote{В~силе
  с~13.08.2025; заменило 375-П. К~12.11.2025 кредитные организации обязаны были
привести ПВК к~новому составу.} задаёт действующий состав ПВК~\cite{cbr-860p}.

ПВК по~860-П состоит из~шести связанных программ. Программа идентификации
клиента, представителя, выгодоприобретателя и~бенефициарного владельца отвечает
на~вопрос, кого впускаем и~какие сведения о~нём фиксируем. Программа управления
риском задаёт классификатор клиентов и~операций по~уровням риска и~пересмотр
этих уровней по~событиям. Программа выявления операций, подлежащих обязательному
контролю, подозрительных операций и~подозрительной деятельности описывает
критерии срабатывания, на~которые опирается мониторинг. Программа
документального фиксирования закрепляет состав сведений, формат и~сроки хранения
для каждого решения. Программа обучения определяет учебный цикл сотрудников
по~ПОД/ФТ. Программа проверки внутреннего контроля задаёт регламент самопроверки
и~регулярного отчёта ответственного сотрудника.

В~коде это разделение проявляется явно. Программа выявления операций живёт
в~модуле классификатора, где сходятся пороги статьи~6, признаки подозрительных
операций статьи~7 и~набор перечней лиц, к~которым применяется обязательный
контроль независимо от~суммы. Журнал аудита (\textit{audit trail}) фиксирует
не~только решение, но~и~стабильный код причины, по~которому это решение
принято~--- иначе разбор возражения клиента, запрос регулятора или
ретроспективная оценка эффективности правил превращается в~раскопки логов:

\codeexcerpt{Python}{samples/ch26-aml/aml_classifier.py}{71}{92}

Разделение статьи~6 и~статьи~7 в~коде уложено в~три ветки: формальный
обязательный контроль по~сумме и~категории операции (ст.~6), формальное лицо
из~перечня (тоже ст.~6) и~подозрительная операция по~правилу из~ПВК (ст.~7).
Каждая ветка кладёт стабильный код причины (\texttt{threshold:cash\_deposit},
\texttt{party\_in\_list}, \texttt{rba:high\_risk\_client}) в~Verdict. Программа
документального фиксирования по~860-П опирается на~эти коды: оператор
не~пересказывает решение словами, а~ссылается на~запись.

ПВК~--- внутренний документ организации, но~Банк России и~Росфинмониторинг
проверяют его как минимум по~двум измерениям: наличие всех шести программ
в~актуальной редакции и~фактическое соответствие правил тому, как работает
система. Несоответствие декларации и~реализации~--- типовое предписание
по~ПОД/ФТ.

\subsection*{Риск-ориентированный подход}

Программа управления риском в~860-П требует от~банка применить
риск-ориентированный подход (\textit{risk-based approach}, RBA) на~всём
периметре ПОД/ФТ. Принцип сформулирован в~Рекомендации~1 FATF и~развёрнут
в~отраслевом руководстве FATF для банковского сектора (октябрь
2014~года)~\cite{fatf-recs,fatf-rba}; в~российской регуляторике RBA закреплён
499-П (идентификация и~оценка риска клиента)~\cite{cbr-499p} и~860-П (применение
оценки к~мониторингу и~факультативному контролю)~\cite{cbr-860p}.

RBA меняет три параметра системы. Глубину сбора сведений на~входе: для низкого
риска оператор обходится минимальным набором, для высокого~--- запрашивает
расширенный пакет, привлекает внешние источники и~оценивает структуру владения.
Чувствительность мониторинга: для высокорискового клиента порог факультативного
контроля сдвигается к~меньшим суммам, а~сами правила срабатывают на~большем
числе сценариев. Периодичность обновления данных и~пересмотра уровня риска: чем
выше уровень, тем короче цикл (см.~раздел~\ref{sec:merchant-monitoring}).

В~примере выше RBA проявляется веткой \texttt{rba:high\_risk\_client}: операция,
формально не~дотягивающая до~порога обязательного контроля, для высокорискового
клиента уходит в~очередь подозрительных. Коэффициент~--- половина, треть,
четверть формального порога~--- задаётся не~законом, а~программой управления
риском. ПВК документирует коэффициент и~обоснование, журнал аудита фиксирует,
что решение принято именно по~этой ветке. На~запрос регулятора банк предъявляет
конкретный коэффициент и~конкретную запись журнала, а~не формулу \enquote{мы
внимательно смотрим за~высокорисковыми клиентами}.

Сведение RBA к~работе с~политически значимыми лицами (PEP, \textit{Politically
Exposed Person}) занижает принцип. PEP~--- один из~триггеров высокого риска
среди прочих; RBA охватывает все сегменты клиентской базы и~все типы операций.

\needspace{8\baselineskip}
\subsection*{Алгоритмы детекции дробления}

Программа выявления операций ловит не~только формальные критерии статьи~6,
но~и~поведенческие паттерны. Самый частый из~них~--- дробление
(\textit{structuring}): клиент разбивает крупную сумму на~несколько операций
ниже порога, чтобы избежать формального обязательного контроля.

Методические рекомендации Банка России от~21.07.2017~№~19-МР
\enquote{О~корпоративных картах юридических лиц и~ИП} прямо обязывают кредитные
организации проводить \enquote{не~реже одного раза в~неделю мониторинг операций
  \dots{} включая анализ характеристик операций, указывающих на~возможное
дробление (разделение) крупных сумм на~меньшие операции}~\cite{cbr-19mr}. Парные
Методические рекомендации от~21.07.2017~№~18-МР \enquote{О~подходах к~управлению
риском ПОД/ФТ} задают общий подход к~RBA в~этой работе~\cite{cbr-18mr};
Методические рекомендации Банка России от~29.02.2024~№~4-МР \enquote{Об~усилении
контроля за~операциями физических лиц} расширяют его на~переводы
физлиц~\cite{cbr-4mr}. Конкретные коды необычных операций, по~которым из~этих
сигналов рождается ФЭС, перечислены в~приложении к~тому же
Приказу~№~18~\cite{rosfinmonitoring-oes}.

Три правила на~скользящих окнах закрывают типовое механическое дробление.

\textbf{Sub-threshold velocity}~--- индекс по~одному счёту-источнику с~подсчётом
числа операций и~общей суммы в~окне. Срабатывает, когда за~окно (обычно неделя
по~19-МР, иногда сутки или час у~активных каналов) накапливается N~операций
общей суммой выше порога обязательного контроля при том, что каждая отдельная
операция меньше порога. Базовый и~самый частый паттерн.

\codeexcerpt{Python}{samples/ch26-aml/structuring.py}{29}{58}

\textbf{Fan-out}~--- один счёт-источник в~коротком окне переводит небольшие
суммы на~много разных счетов. Часто видна как смурфинг (\textit{smurfing}) через
подставные счета: один реальный плательщик расщепляет поток через сеть знакомых
или контролируемых аккаунтов.

\codeexcerpt{Python}{samples/ch26-aml/structuring.py}{61}{87}

\textbf{Fan-in}~--- зеркальный паттерн: много разных счетов-источников переводят
небольшие суммы на~один счёт-получатель. Та же логика смурфинга
с~противоположной стороны: получатель собирает деньги через раздробленный фронт.

\codeexcerpt{Python}{samples/ch26-aml/structuring.py}{90}{116}

Каждое правило держит индекс по~своему ключу: \texttt{SubThresholdVelocity}
и~\texttt{FanOut} индексируют по~\texttt{source}, \texttt{FanIn}~---
по~\texttt{destination}. Окна не~зависят друг от~друга, и~срабатывание одного
правила не~мешает другому: один и~тот же поток может одновременно поднять
\texttt{sub\_threshold} (по~сумме за~неделю) и \texttt{fan\_out} (по~числу
получателей за~тот же интервал).

\textbf{Графовые модели идут дальше}. Метод GARG-AML (\textit{Graph-Aided Risk
Grading for Anti-Money Laundering}, 2025) строит над матрицей смежности
транзакционного графа признаки плотности (\textit{density-based}) и~распознаёт
восемь типов структур: fan-out, fan-in, gather-scatter, scatter-gather, simple
cycle, bipartite, stack и random~\cite{garg-aml-2025}. Принципиальная
разница~--- глубина: правила выше ловят дробление на~одном переходе (источник
раздаёт получателям или собирает от~них), GARG-AML видит цепочки от~A через
посредников к~B, где каждый отдельный шаг подпороговый, но~связность графа
выдаёт намерение. Платится за~это стоимостью поддержания актуального графа
платежей и~снижением интерпретируемости сигналов. У~большинства банков графовая
(\textit{graph-based}) надстройка работает поверх правил (\textit{rule-based}),
не~заменяя их.

\highlight{Каждый сработавший паттерн уходит в~ту же очередь ручного разбора,
что и~формальные сигналы по~ст.~6} (раздел~\ref{sec:aml-false-positives}).
Алгоритм поднимает запись с~кодом причины (\texttt{sub\_threshold},
\texttt{fan\_out}, \texttt{fan\_in}) и~аналитик второй линии решает, заводить
ФЭС или закрывать как FP. Без RBA-калибровки доля FP по~дроблению особенно
высока: интенсивность, которая для одного клиента~--- нормальный поток, для
другого~--- явная схема обхода порога.

Банк России при~надзоре за~оператором платёжной инфраструктуры проверяет рабочую
систему ПОД/ФТ целиком: формальные правила, обученных сотрудников и~маршрут
обработки случаев.

\needspace{8\baselineskip}
\section{Ложные срабатывания финмониторинга}\label{sec:aml-false-positives}

Программа выявления операций срабатывает чаще, чем нужно. Доля ложных
срабатываний (\textit{false positives}, FP) в~финмониторинге у~банков
по~отраслевым обсуждениям держится в~районе 85--95\,\% от~общего числа поднятых
сигналов, а~в~системах на~одних правилах без~RBA-калибровки поднимается
до~95--98\,\%; конкретные цифры оператор не~раскрывает, порядок виден по~отчётам
вендоров и~консультантов (Nice Actimize, FICO TONBELLER, SAS, Oracle FCCM).
Низкая базовая частота настоящего отмывания делает поток FP штатным выходом
любой разумно чувствительной программы выявления.

\highlight{Маршрут разбора отделяет сигнал от~транзакции}. Сигнал поднимается
классификатором, попадает в~очередь аналитика первой линии, аналитик собирает
контекст (профиль клиента, история операций, связанные стороны, внешние
источники) и~принимает одно из~трёх решений. Закрыть как FP, оставив запись
в~журнале с~кодом причины. Передать на~вторую линию для углублённой проверки
(\textit{enhanced due diligence}, EDD). Принять решение о~направлении ФЭС либо
о~применении меры по~115-ФЗ. Сама операция в~это время либо стоит в~приостановке
(если правило приостанавливает зачисление), либо уходит дальше, а~разбор идёт
по~записи журнала.

SLA задаётся характером сигнала. Подозрение на~финансирование терроризма~---
немедленный разбор и~доклад ответственному сотруднику. Стандартный сигнал
по~правилу мониторинга~--- 1--3 рабочих дня. Периодический пересмотр уровня
риска клиента~--- по~графику программы управления риском. Сроки нужны
под~операционный риск: операция, по~которой решение задерживается, либо
накапливает риск клиентской жалобы (если стоит в~приостановке), либо успевает
уйти в~зачисление и~блокировка перестаёт быть доступной.

Обработка одного сигнала на~первой линии~--- порядка одного-двух человеко-часов,
на~второй~--- от~половины рабочего дня и~больше. На~потоке в~10\,000~сигналов
в~сутки это десятки FTE только под~AML-разбор; в~крупном банке число выражается
сотнями. Каждый процентный пункт FP-уровня переводится в~FTE напрямую.

Тихое снижение чувствительности~--- единственный быстрый способ выровнять
очередь, и~единственный способ за~квартал довести систему до~регуляторного
нарушения. Сдвиг порога вверх без обновления риск-оценки сокращает FP
пропорционально, но~и сокращает обнаружение настоящих случаев. Регулятор
проверяет программу выявления по~эффективности, а~не по~числу разобранных
сигналов~--- сколько направленных ФЭС привели к~делам, сколько мер
по~замораживанию применено по~результатам собственного мониторинга. Программа,
в~которой за~квартал не~направлено ни одного ФЭС, выглядит подозрительно ровно
в~той же мере, в~какой подозрительна программа, направившая сто тысяч.

Устойчивый способ снизить FP~--- добавлять признаки в~оценку риска и~менять
структуру правил под~актуальные сценарии (типологии 860-П, информационные письма
Росфинмониторинга, опубликованные расследования ЕАГ). Чем точнее RBA-калибровка,
тем меньше шума и~тем больше попаданий в~настоящие случаи. Это задача программы
управления риском, а~не оператора первой линии: первая линия не~имеет права
менять пороги, она имеет право только закрыть сигнал и~зафиксировать причину.

\practicenote{%
  Очередь AML-разбора измеряется в~часах и~FTE, а~не в~числе сигналов. Если
  очередь растёт быстрее найма аналитиков, у~команды две недели, чтобы
  пересмотреть RBA-калибровку либо смириться с~ростом отказов по~необработанным
  сигналам.

}
\section{Отзыв лицензии у~КИВИ~Банка}\label{sec:aml-qiwi}

21~февраля 2024~года Банк России отозвал лицензию у~КИВИ~Банка Приказом
от~21.02.2024~№~ОД-266~\cite{cbr-qiwi-license-2024}. В~официальном пресс-релизе
регулятор указал на~систематические нарушения требований законодательства
по~ПОД/ФТ, вовлечённость банка в~высокорисковые операции, включая переводы
в~пользу криптообменников и~нелегальных онлайн-казино, а~также на~случаи
открытия кошельков QIWI с~использованием персональных данных граждан без их
ведома.

115-ФЗ требует от~банков идентифицировать клиентов и~проводить контроль
операций; открытие кошельков без ведома граждан означает, что процедура
идентификации либо не~выполнялась, либо не~проверялась системой контроля.
Переводы в~пользу криптообменников и~казино становятся проблемой тогда, когда
мониторинг транзакций (\textit{transaction monitoring}~--- автоматическая
проверка операций против набора правил и~рисковых моделей) не~маркирует их как
высокорисковые или маркирует, но~событие не~попадает на~ручной разбор. Оба сбоя
технические: отсутствие автоматического триггера для усиленной проверки клиента
(\textit{enhanced due diligence}, EDD) и~неполный журнал аудита (упорядоченная
  последовательность записей, позволяющая восстановить историю принятых решений
по~транзакции) по~спорным транзакциям.

\importantnote{Законодательный ответ на~ту~же практику~--- Федеральный
  закон от~20.02.2026~№~44-ФЗ \enquote{Об~ограничении идентификации физических лиц
  банковскими платёжными агентами}~\cite{fz-44}: с~03.03.2026 БПА не~вправе
  идентифицировать физлиц, выдавать карты и~открывать электронные кошельки.
  Пояснительная записка опирается на~выявленные Банком России случаи незаконного
  использования персональных данных для~открытия кошельков без~ведома граждан
  и~на~борьбу с~дропперством. Идентификация физлица возвращена внутрь периметра
  банка; БПА-агрегатор по~ст.~14.1 161-ФЗ сохраняет идентификацию
ЮЛ/ИП-субмерчантов.}

Операционный риск ПОД/ФТ для платёжного бизнеса выходит за~штрафы: отзыв
лицензии у~ключевого расчётного банка одномоментно обнуляет каналы приёма
и~выплат у~всех его клиентов. Бизнесу с~зависимостью от~одного платёжного
провайдера нужен заранее проверенный план непрерывности~--- альтернативный
провайдер, резервные счета и~маршрут переключения.

\section{FATF и~ЕАГ для СНГ}\label{sec:aml-fatf-eag}

FATF задаёт международные стандарты через рекомендации, взаимные оценки
и~публичные списки юрисдикций под~усиленным
наблюдением~\cite{fatf-recs,fatf-increased-monitoring}. С~24~февраля 2023~года
членство Российской Федерации в~FATF
приостановлено~\cite{fatf-russia-suspension-2023}: формально это означает, что
Россия не~участвует в~выработке решений FATF, не~направляет своих представителей
в~её рабочие группы и~не подвергается очередным взаимным оценкам в~стандартном
порядке. В~СНГ роль региональной \textit{FATF-style body} (организации,
действующей по~методологии FATF) играет ЕАГ~--- Евразийская группа
по~противодействию легализации преступных доходов и~финансированию терроризма.
ЕАГ проводит взаимные оценки государств региона и~публикует методологические
документы, на~которые опираются национальные регуляторы стран
СНГ~\cite{eag-overview,eag-mutual-evaluations}.

После приостановки членства стандарты FATF продолжают влиять на~платёжный бизнес
в~России косвенно. Банки-корреспонденты и~платёжные посредники в~третьих
юрисдикциях смотрят на~риск через рекомендации FATF и~публичные списки. Крупные
банки практикуют политику снижения контрагентского риска
(\enquote{\textit{de-risking}}~--- отказ от~целых классов клиентов или
юрисдикций вместо точечного управления риском по~каждому отношению) и~сокращают
или полностью прекращают отношения с~целыми классами клиентов или
юрисдикциями~\cite{fatf-correspondent-banking}. Для российского оператора это
означает усложнение открытия и~поддержания корреспондентских счетов в~валютах
недружественных стран~--- дополнительные проверки на~стороне контрагента
и~снижение доступности канала.

\practicenote{%
  Перед выходом в~новую юрисдикцию проверяются оценка ЕАГ для стран СНГ либо
  FATF для остальных, последняя публичная взаимная оценка страны и~местная
  практика ПОД/ФТ. Если что-то из~этого вызывает сомнение, заранее закладывается
  запас по~времени и~каналам расчёта.

}

\section{AML и~блокчейн-аналитика как готовый сервис}\label{sec:aml-industry-providers}

Мониторинг транзакций по~115-ФЗ можно реализовать своими силами, но~крупный
сегмент рынка покупает его готовым процессом. Платформы класса
ComplyAdvantage~\cite{complyadvantage-overview} объединяют в~одном API правила
обязательного контроля, RBA-скоринг, санкционный скрининг (\textit{sanctions
screening}) и~поиск негативных упоминаний (\textit{adverse media}). Составление
и~поддержка ПВК остаются за~продуктом: по~букве 115-ФЗ они не~могут быть
делегированы; технические части мониторинга и~актуализация санкционных списков
уходят наружу.

Крипто-сегмент закрывает отдельный класс инструментов~--- блокчейн-аналитика.
Chainalysis~\cite{chainalysis-overview}~--- крупнейший игрок по~охвату
и~использованию в~правоохранительных расследованиях; рядом работают Elliptic
(на~рынке с~2013~года, с~декларируемым покрытием около 99~\% активов)
и~TRM~Labs, выросший с~2018~года до~оценки уровня \textit{unicorn} и~сделавший
ставку на~ИИ-скоринг и~межсетевой трейс в~реальном времени. Эти инструменты
отслеживают происхождение средств по~истории в~блокчейне (\textit{on-chain}),
маркируют адреса миксеров (\textit{mixer}~--- сервис анонимизации
криптоактивов), санкционных кошельков и~нелегальных сервисов и~подмешивают этот
сигнал в~авторизационный или послеавторизационный мониторинг.

Выбор \enquote{строить или~купить} в~AML-мониторинге решается по~роли функции
в~бизнесе. Компании, для которых AML~--- сопутствующая функция, почти всегда
выбирают готовый сервис; банки и~криптобиржи, для которых качество AML-аппарата
прямо определяет лицензию, чаще держат гибрид~--- внешние списки и~трейсинг
внутри собственного движка правил.

Выбор задаёт состав функций, которые уходят к~внешнему поставщику; размер банка
определяет сумму договора, но~не предмет. Готовый сервис закрывает санкционный
скрининг, поиск негативных упоминаний и~базовую RBA-разметку клиента. ПВК,
журнал аудита и~решение о~направлении ФЭС остаются за~оператором по~115-ФЗ
и~860-П~--- делегировать их нельзя по~букве закона независимо от~того, насколько
удобный API предлагает вендор. Гибрид нужен там, где собственные правила
выявления зависят от~продукта: переводы между картами, корпоративные счета,
обмен криптоактивов, маркетплейс с~субмерчантами. Для каждого из~этих сценариев
готовый сервис даёт правила и~списки, а~связки с~платёжным журналом и~порядок
сборки ФЭС всё равно собираются внутри.

Блокчейн-аналитика отвечает на~вопрос, на~который у~банка нет источника
по~115-ФЗ: связан~ли криптокошелёк, с~которого пришли средства, с~миксером,
даркнет-площадкой, выкупом вымогательскому ПО (\textit{ransomware}) или
санкционным адресом. Chainalysis и~Elliptic ведут таблицу связей адресов
по~собственной кластеризации и~открытым типологиям; банк получает оценку риска
кошелька и~метки источника, не~раскрывая отдельные операции. Когда клиент
пополняет счёт переводом с~криптобиржи или выводит на~неё средства, эта оценка
становится отдельным входом в~AML-скоринг банка наравне с~типологией клиента
и~MCC контрагента-криптобиржи. Чем выше доля переводов между фиатным счётом
клиента и~криптобиржей в~портфеле, тем чаще решение о~направлении ФЭС зависит
именно от~этой оценки.
